\chapter{Konzeption und Methodik}
\label{chap:Konzeption und Methodik}

\section{Python}

Python ist eine weitverbreitete Programmiersprache, die von der Python Software Foundation weiterentwickelt wird. Ihr Designer ist Guido van Rossum. Die Sprache wurde in den 90er-Jahren entwickelt und hätte ursprünglich als Skriptsprache für das Betriebssystem Amoeba gedient. 1994 erschien Python 1.0. Die aktuellste Version ist 3.14.8. Python eignet sich für Anfänger und für Experten, weil die Sprache leicht verständlich ist und viele Libraries (Bibliotheken) hat, aber auch komplexe Aufgaben lösen kann. In Python ist jeder Datentyp ein Objekt.

\section{Libraries}

\subsection{Socket}

Socket ist eine Bibliothek, welche Verbindungen erstellt mit Server und Client. Sockets kommen ursprünglich aus dem ARPANET (Advanced Research Projects Agency Network), wo schon seit 1971 benutzt wurden. Sockets arbeiten mit dem TCP/IP, weswegen sie eine zuverlässige Verbindung zischen mehreren Geräten erzeugen können. Das Socket übernimmt alle tiefen Schichten des OSI-Models, von der Sicherungsschich bis hin zur Sitzungsschicht.

\subsection{Threading}

Threading ist eine Bibliothek, welche erlaubt mehrere Operationen Scheinweise parallel zu bearbeiten. Es ermöglicht Prozesse die einen Stau bei der Ausführung verursachen könnten zu verhindern. Indem es diese als separate Operationen aufteilt.

\subsection{Tkinter}

Tkinter ist eine Bibliothek, welche eine ressourcenschonendes GUI (graphical user interface) erstellt, die vollkommen selbst gestaltbar ist. Mit dieser kann man Fenster bauen, die alle benötigten Nutzereingaben nehmen und weiterleiten können. Im OSI-Model wäre sie die Anwendungsschicht. Es gibt auch eine erweiterte Bibliotheken, die Tkinter als Fundament verwendet, aber zusätzlich die GUI attraktiver gestalten. Beispielsweise das CustomTkinter, welches in dieser Arbeit verwendet wird.

\subsection{Pycryptodomex}

Pycryptodomex ist eine Bibliothek, welche eine Vielzahl an modernen Verschlüsselungen beinhaltet. Sie ist aber ein Fork vom alten Pycryptodome, weil es oft zu Bibliothekkonflikten mit dem original kommt. Diese Bibliothek beinhaltet alle Algorithmen, welche für jede Verschlüsselung benötigt wird.